\section{Conclusion}

\subsection{Summary}

We presented NEXUS, a modular framework for combining six sparse attention 
mechanisms. NEXUS achieves sub-quadratic attention for encoders while 
falling back to dense for causal models.

\subsection{Results Recap}

\begin{table}[H]
\centering
\small
\begin{tabular}{lll}
\toprule
\textbf{Model} & \textbf{Acc $\Delta$} & \textbf{Path} \\
\midrule
DistilBERT & -39.79\% & sparse \\
RoBERTa & -27.64\% & sparse \\
GPT-2 & -27.75\% & dense \\
\bottomrule
\end{tabular}
\end{table}

\subsection{Key Takeaways}

\begin{enumerate}
\item \textbf{Modular design enables experimentation.}
\item \textbf{Encoder-only sparsity} in our implementation.
\item \textbf{Two components (stripe, RFF) inactive} under default SST-2.
\item \textbf{Softpick needs care.}
\item \textbf{Speedup grows with $N$.}

\end{enumerate}
\subsection{Future Work}

\subsubsection{Short-Term}
\begin{itemize}
\item Causal block-sparse path
\item Enable stripe for shorter sequences
\item Alternative to softpick

\end{itemize}
\subsubsection{Long-Term}
\begin{itemize}
\item Component interaction study
\item Theoretical analysis
\item Multi-task evaluation

\end{itemize}
\subsection{Closing Remarks}

NEXUS is a research framework for studying combined sparse attention. We 
release code and checkpoints, and we explicitly document which components 
are active in our experiments.

\subsection{Acknowledgments}

We thank the anonymous reviewers and the compute cluster administrators for their support.

\subsection{Author Contributions}

\begin{itemize}
\item \textbf{Taha Majs} designed, implemented, ran experiments, wrote paper.
\item \textbf{Collaborative Research Team} provided feedback, edited.

\end{itemize}
\subsection{Conflict of Interest}

None.
\begin{lstlisting}

---
\end{lstlisting}